\chapter{Realización constitutiva del vacío y propagación electromagnética–electrodébil}
\label{chap:parte-iv-52}

La estructura bidireccional de la acción induce relaciones constitutivas del vacío y separa los sectores de propagación, impedancia, carga y mezcla electrodébil. Las referencias SI aparecen como realización posterior de funciones internas ya determinadas.

\section{Realizaci\'on electromagn\'etica del operador bidireccional de acci\'on: \(\varepsilon\), \(\mu\), \(Z\) y \(c\)}
\label{rm:ch:vacuum}

\subsection{Derivaci\'on de la escala de acci\'on a partir del estado dodecaf\'asico}

La d\'ecada de acci\'on se determina intr\'insecamente, sin adoptarla como
unidad heredada. El funcional determinantal local se aplica a la salida HMT
previamente construida \(\alpha_{\HMT}\) y define
\begin{equation}
\Sigma_H=\varphi_{\HMT}\alpha_{\HMT}^{16}.
\label{rm:eq:action-determinantal-reader}
\end{equation}
El exponente \(16\) es el orden del cokernel local del cierre firmado; su valoraci\'on decimal fija
\begin{equation}
\operatorname{ord}_{10}(\Sigma_H)=-34.
\label{rm:eq:action-decade}
\end{equation}

El origen del exponente se puede verificar dentro de este volumen. Sobre la
\'orbita local de cuatro hojas act\'ua el bloque de Hadamard
\begin{equation}
H_4=
\begin{pmatrix}
1&1&1&1\\
1&1&-1&-1\\
1&-1&1&-1\\
1&-1&-1&1
\end{pmatrix}.
\label{rm:eq:action-hadamard}
\end{equation}

\begin{lemma}[\'Indice local de la recta de acci\'on]
La forma normal de Smith de \(H_4\) es
\begin{equation}
\operatorname{SNF}(H_4)=\operatorname{diag}(1,2,2,4).
\label{rm:eq:action-smith}
\end{equation}
En consecuencia,
\[
\operatorname{coker}(H_4)\simeq
\Z/2\Z\oplus\Z/2\Z\oplus\Z/4\Z,
\qquad
|\operatorname{coker}(H_4)|=16.
\]
\end{lemma}

\begin{proof}
El m\'aximo com\'un divisor de las entradas es uno; el de los menores
\(2\times2\) no nulos es dos; el de los menores \(3\times3\) es cuatro; y
\(|\det H_4|=16\). Los cocientes sucesivos de estos divisores determinantes
son \(1,2,2,4\), que prueban \cref{rm:eq:action-smith}. El producto de los
invariantes da el orden del con\'ucleo.
\end{proof}

La norma de la secci\'on \(\alpha_{\HMT}\) sobre esas diecis\'eis hojas
produce \(\alpha_{\HMT}^{16}\); la autoescala \(\varphi_{\HMT}\) act\'ua
una sola vez sobre la base, no una vez por hoja. De ah\'i
\cref{rm:eq:action-determinantal-reader}. Las comparaciones certificadas
\begin{equation}
\alpha_{\HMT}^{16}<10^{-34}
<\varphi_{\HMT}\alpha_{\HMT}^{16}<10^{-33}
\label{rm:eq:action-decade-bounds}
\end{equation}
demuestran, sin introducir la unidad SI, que el orden decimal es
exactamente \(-34\). La evaluaci\'on terminal comienza por
\begin{equation}
\Sigma_H
=1.046243838104756580184229208303089\ldots\times10^{-34}.
\label{rm:eq:action-sigma-value}
\end{equation}
La sustituci\'on del con\'ucleo por tres copias con exponente \(16^3\), o la
aplicaci\'on reiterada de \(\varphi\) en cada hoja, modificar\'ia la
genealog\'ia. Ambas operaciones constituyen criterios de refutaci\'on y no
convenciones equivalentes.
La cadena causal que se conserva es, por tanto,
\[
\APP\to\TRIT\to\TPK\to U_{12}^{\rm sign}
\to\alpha_{\HMT}\to\Sigma_H
\to10^{-34}\U_S\to
\{\hbar_{\rm pre},\hbar_{\rm ret}\}.
\]
Este volumen no rederiva \(\alpha_{\HMT}\), sino que desarrolla la geometr\'ia
metrol\'ogica posterior a la selecci\'on del tipo y de la d\'ecada por
\(\Sigma_H\). El entero \(-34\) no es una mantisa ni un ajuste al Sistema
Internacional, sino la valoraci\'on terminal del funcional local.

\subsection{Secciones orientadas de la recta de acci\'on}

Las realizaciones anterior y posterior al retorno constituyen dos secciones
de una misma recta de acci\'on:

\begin{equation}
\hbar_{\rm pre}=H_5\,10^{-34}\U_S,\qquad
\hbar_{\rm ret}=\eta_{\rm ret}\,10^{-34}\U_S.
\label{rm:eq:action-twoway}
\end{equation}

La d\'ecada com\'un determina la recta, mientras que la informaci\'on
orientada se expresa mediante

\begin{equation}
R_{\rm act}=\frac{\hbar_{\rm pre}}{\hbar_{\rm ret}}
=\frac{H_5}{\eta_{\rm ret}},
\qquad
\xi_{\rm act}=\frac12\log R_{\rm act}.
\label{rm:eq:action-rapidity}
\end{equation}

La primera cantidad proporciona una coordenada multiplicativa y la segunda
una coordenada aditiva de tipo rapidez. No representan dos constantes de
Planck independientes, sino las dos orientaciones de una misma secci\'on
respecto del retorno.

Si \(h_a=2\pi\hbar_a\), \(a\in\{\mathrm{pre},\mathrm{ret}\}\), una magnitud proporcional a \(h_a^{1/2}\) cambia como \(R_{\rm act}^{1/2}\), una proporcional a \(h_a^{-1/2}\) cambia como \(R_{\rm act}^{-1/2}\), y un cociente donde la acci\'on se cancela es bidireccional invariante.

\subsection{Descomposici\'on espectral bicapa del operador angular}

Sea \(R=R^*=R^{-1}\) la involuci\'on activa y

\[
P_\pm=\frac{I\pm R}{2},
\qquad
T=q_+P_+ +q_-P_-,
\qquad 0<q_+<q_-<1.
\]

Con
\[
x=A_{\rm rad}=\frac{\pi A}{180},\qquad
y=C^*_{\rm rad}=\frac{\pi C^*}{180},
\]
los dos canales se construyen antes de evaluar el vac\'io:
\begin{equation}
q_+=e^{-x-y},\qquad q_-=e^{-x+y}.
\label{rm:eq:vacuum-angular-channels}
\end{equation}
El intercambio de hojas invierte \(C^*\), es decir \(y\mapsto-y\), y permuta \(q_+\leftrightarrow q_-\). La pareja es as\'i la lectura espectral de un mismo antecedente angular, no dos par\'ametros ajustados por separado.

\begin{definition}[Respuesta constitutiva]
\begin{equation}
\mathcal V_{90,120}=(I-T^{90})(I-T^{120})^{-1}.
\label{rm:eq:vacuum-operator}
\end{equation}
\end{definition}

Como \(\|T\|<1\), el denominador es positivo e invertible. El c\'alculo funcional da

\begin{equation}
\mathcal V_{90,120}=r_+P_+ +r_-P_-,
\qquad
r_\pm=\frac{1-q_\pm^{90}}{1-q_\pm^{120}}.
\label{rm:eq:vacuum-eigenvalues}
\end{equation}

La afirmaci\'on de invertibilidad merece hacerse expl\'icita. Para cada hoja,
\(0<q_\pm^{120}<1\); por tanto
\[
I-T^{120}=(1-q_+^{120})P_++(1-q_-^{120})P_-
\]
es estrictamente positivo y
\[
(I-T^{120})^{-1}
=\frac{P_+}{1-q_+^{120}}+
 \frac{P_-}{1-q_-^{120}}.
\]
La respuesta \(\mathcal V_{90,120}\) no se obtiene resolviendo cuatro
ecuaciones escalares separadas: es el resultado del c\'alculo funcional de una
sola contracci\'on positiva. Esta observaci\'on impide que permeabilidad,
permitividad, impedancia y propagaci\'on pierdan el antecedente que comparten.

\begin{theorem}[Determinaci\'on espectral conjunta de las cuatro relaciones constitutivas]
Los cuatro caracteres constitutivos
\begin{equation}
\widehat\varepsilon=r_+^2,\qquad
\widehat\mu=r_-^2,\qquad
\widehat Z=\frac{r_-}{r_+},\qquad
\widehat c=\frac1{r_+r_-}
\label{rm:eq:vacuum-four}
\end{equation}
son funciones espectrales de \(\mathcal V_{90,120}\). Satisfacen
\begin{equation}
\widehat\mu\widehat\varepsilon=\widehat c^{-2},
\qquad
\frac{\widehat\mu}{\widehat\varepsilon}=\widehat Z^2.
\label{rm:eq:vacuum-relations}
\end{equation}
\end{theorem}

\begin{proof}
Las identidades se obtienen sustituyendo \cref{rm:eq:vacuum-four}. No interviene ninguna carta dimensional ni ning\'un valor exterior.
\end{proof}

\begin{corollary}[Reconstrucci\'on de las dos hojas]
Los cuatro caracteres no contienen cuatro grados de libertad. Cualquiera de
los pares \((\widehat\mu,\widehat\varepsilon)\),
\((\widehat Z,\widehat c)\) o \((\mathcal E,\mathcal B)\) reconstruye los
dos autovalores positivos:
\begin{align}
r_-&=\sqrt{\widehat\mu}
=\sqrt{\widehat Z/\widehat c},
&r_+&=\sqrt{\widehat\varepsilon}
=\frac1{\sqrt{\widehat Z\widehat c}},
\label{rm:eq:vacuum-reconstruct-r}\\
\log r_-&=\frac{\mathcal E+\mathcal B}{2},
&\log r_+&=\frac{\mathcal E-\mathcal B}{2}.
\label{rm:eq:vacuum-reconstruct-log}
\end{align}
En particular, la orientaci\'on se pierde si se conserva s\'olo el producto,
pero se recupera al a\~nadir el cociente.
\end{corollary}

\begin{proof}
La primera l\'inea es la ra\'iz positiva de las identidades
\cref{rm:eq:vacuum-four,rm:eq:vacuum-relations}. La segunda invierte el sistema
lineal
\(\mathcal E=\log r_-+\log r_+\),
\(\mathcal B=\log r_--\log r_+\).
\end{proof}

Este corolario expresa una propiedad central de HMT: producto y cociente son
funcionales complementarios de un mismo estado. El producto determina la
propagaci\'on com\'un, mientras que el cociente conserva la orientaci\'on
relativa de las hojas interior y exterior.

En el certificado V2,

\begin{align*}
r_+&=0.9999996285482493631786952281\ldots,\\
r_-&=0.9997241371895183641315165853\ldots,\\
\widehat\varepsilon&=0.9999992570966367027604416155\ldots,\\
\widehat\mu&=0.9994483504793269350899815559\ldots,\\
\widehat Z&=0.9997245085389372154555543466\ldots,\\
\widehat c&=1.0002763104861575261063862689\ldots.
\end{align*}

Estos numerales son el reconocimiento terminal de \cref{rm:eq:vacuum-operator}; no definen sus exponentes \(90,120\).

\subsection{Factorizaci\'on arm\'onica de los sectores \texorpdfstring{\(90\)}{90} y \texorpdfstring{\(120\)}{120}}

El semigrupo arm\'onico es

\[
\langle90,120\rangle=30\langle3,4\rangle.
\]

Si \(s=q^{30}\), entonces

\begin{equation}
\frac{1-q^{90}}{1-q^{120}}
=\frac{1-s^3}{1-s^4}
=\frac{1+s+s^2}{1+s+s^2+s^3}.
\label{rm:eq:three-four-completion}
\end{equation}

La identidad tambi\'en separa la respuesta de su defecto. Si
\[
F_{3\to4}(s)=\frac{1+s+s^2}{1+s+s^2+s^3},
\qquad 0<s<1,
\]
entonces
\begin{equation}
d_{3\to4}(s):=1-F_{3\to4}(s)
=\frac{s^3}{1+s+s^2+s^3}>0.
\label{rm:eq:vacuum-defect34}
\end{equation}
El t\'ermino adicional se conserva exactamente como defecto
positivo. Adem\'as,
\begin{equation}
F_{3\to4}'(s)
=-\frac{s^2(3+2s+s^2)}{(1+s+s^2+s^3)^2}<0.
\label{rm:eq:vacuum-monotonic34}
\end{equation}
Por consiguiente, el orden de los canales angulares determina un\'ivocamente
el orden de las respuestas constitutivas. La asignaci\'on de
\(\widehat\mu\) y \(\widehat\varepsilon\) a las hojas respectivas se obtiene
anal\'iticamente, sin recurrir a sus expansiones decimales.

Cada hoja del vac\'io constituye la completaci\'on exacta de tres a cuatro
t\'erminos de un mismo modo elemental \(30\). En consecuencia, permeabilidad
y permitividad se interpretan como dos funcionales cuadr\'aticos de una misma
completaci\'on evaluada sobre hojas conjugadas, y no como par\'ametros
constitutivos independientes.

Como la funci\'on \((1-q^{90})/(1-q^{120})\) decrece en \(0<q<1\) y \(q_+<q_-\), se obtiene

\[
0<r_-<r_+<1,\qquad
0<\widehat\mu<\widehat\varepsilon<1,\qquad
0<\widehat Z<1<\widehat c.
\]

Los l\'imites del funcional determinan su geometr\'ia. Cuando \(q\to0^+\),
\(F(q)\to1\): las cuatro lecturas se acercan a la unidad y el defecto de
memoria desaparece. Cuando \(q\to1^-\),
\[
F(q)\longrightarrow\frac{90}{120}=\frac34.
\]
La respuesta toma, por tanto, valores en el intervalo positivo \((3/4,1)\). El
extremo \(3/4\) no es una constante ajustada: es el cociente de las dos
longitudes de canal y la realizaci\'on continua de la completaci\'on
tres--a--cuatro.

\subsection{Involuci\'on de hojas y conservaci\'on de la orientaci\'on}

El intercambio de hojas \(C^*\mapsto-C^*\) act\'ua por

\begin{equation}
\widehat\mu\longleftrightarrow\widehat\varepsilon,\qquad
\widehat Z\longmapsto\widehat Z^{-1},
\qquad
\widehat c\longmapsto\widehat c.
\label{rm:eq:vacuum-sheet}
\end{equation}

La propagaci\'on com\'un es par; la impedancia conserva la orientaci\'on.
Esta distinci\'on reaparecer\'a en Barbero y CKM: no todos los funcionales
asociados al cambio de hoja tienen la misma paridad.

\subsection{Caracteres logar\'itmicos y compatibilidad tracial}

Definamos

\[
\mathcal E=\log(r_+r_-),\qquad
\mathcal B=\log(r_-/r_+).
\]

Con la traza normalizada \(\tau=\Tr/2\),

\begin{equation}
2\tau(\log\mathcal V)=\mathcal E,
\qquad
-2\tau(R\log\mathcal V)=\mathcal B.
\label{rm:eq:vacuum-traces}
\end{equation}

La pareja \((\mathcal E,\mathcal B)\) constituye un sistema de coordenadas lineal del logaritmo
operatorio:
\begin{equation}
\log\mathcal V
=\frac{\mathcal E}{2}I-\frac{\mathcal B}{2}R.
\label{rm:eq:vacuum-log-decomposition}
\end{equation}
La componente escalar es invariante bajo el intercambio de hojas; la
componente orientada cambia de signo. Esta descomposici\'on demuestra
directamente la paridad de la propagaci\'on y el car\'acter orientado de la
impedancia.

La amplificaci\'on non\'adica

\[
\mathcal V_m=\mathcal V\otimes I_{9^m}
\]

es compatible con las expectativas \(E_m=\mathrm{id}\otimes\tau_9\):

\[
E_m(\mathcal V_{m+1})=\mathcal V_m.
\]

Por tanto, \(\mathcal E\) y \(\mathcal B\) no son peculiaridades de una matriz \(2\times2\); permanecen en la torre UHF y en su cierre tracial \(II_1\).

M\'as precisamente, para todo polinomio \(p\) y despu\'es, por continuidad,
para toda funci\'on continua en el espectro,
\begin{equation}
E_m\bigl(p(\mathcal V_{m+1})\bigr)=p(\mathcal V_m).
\label{rm:eq:vacuum-functional-refinement}
\end{equation}
La compatibilidad no preserva solamente dos n\'umeros: preserva todo el
\'algebra funcional generada por la respuesta. En particular conserva sus
proyectores espectrales, sus potencias, su logaritmo y las formas cuadr\'aticas
que miden los defectos \cref{rm:eq:vacuum-defect34}.

\subsection{Realizaci\'on dimensional de las relaciones constitutivas}

La Mec\'anica Dimensional determina

\begin{align}
Z_{\rm vac}&=\widehat Z\,u_Su_Q^{-2},\\
c_{\rm vac}&=\widehat c\,u_C,\\
\mu_{\rm vac}&=\widehat\mu\,u_Su_C^{-1}u_Q^{-2},\\
\varepsilon_{\rm vac}&=\widehat\varepsilon\,u_S^{-1}u_C^{-1}u_Q^2.
\label{rm:eq:vacuum-sections}
\end{align}

Las identidades constitutivas se mantienen en las rectas:

\begin{equation}
\mu_{\rm vac}\varepsilon_{\rm vac}=c_{\rm vac}^{-2},
\qquad
\mu_{\rm vac}/\varepsilon_{\rm vac}=Z_{\rm vac}^2.
\end{equation}

La carta SI es posterior. No es necesario declarar \(c\) primitivo y resolver despu\'es \(\mu\varepsilon=c^{-2}\): los tres objetos descienden conjuntamente de \(\mathcal V\) y de las rectas declaradas.

La separaci\'on de niveles se expresa mediante el diagrama
\[
\begin{array}{c}
(A,C^*)\longrightarrow T\longrightarrow\mathcal V
\longrightarrow
(\widehat\mu,\widehat\varepsilon,\widehat Z,\widehat c)
\\[2mm]
\Big\downarrow\text{ cartas de recta}\hspace{36mm}
\Big\downarrow\text{ evaluaci\'on final}
\\[2mm]
(u_S,u_Q,u_C)\longrightarrow
(\mu_{\rm vac},\varepsilon_{\rm vac},Z_{\rm vac},c_{\rm vac})
\longrightarrow\text{ coordenadas SI}.
\end{array}
\]
La primera l\'inea es adimensional y operatoria; la segunda determina tipos de
magnitud; la tercera selecciona una carta metrol\'ogica. Alterar la carta final
no modifica ni los autovalores ni las identidades constitutivas. En ese
sentido preciso, la unificaci\'on matem\'atica--f\'isica consiste en distinguir
el invariante geneal\'ogico de su sistema de coordenadas, y no en suprimir la
estructura dimensional.

\begin{proposition}[Minimalidad espectral]
Sea \(W\) un operador positivo de dos hojas con espectro simple
\(\{r_+,r_-\}\). Si cuatro caracteres positivos satisfacen
\(\mu\varepsilon=c^{-2}\) y \(\mu/\varepsilon=Z^2\), si la propagaci\'on
com\'un se determina mediante \(c^{-1}=\det W=r_-r_+\), y si su orientaci\'on
fija \(Z=r_-/r_+\), entonces necesariamente
\[
(\mu,\varepsilon,Z,c)
=\left(r_-^2,r_+^2,\frac{r_-}{r_+},\frac1{r_-r_+}\right).
\]
Por tanto \cref{rm:eq:vacuum-four} no es una parametrizaci\'on entre muchas:
es la realizaci\'on positiva m\'inima compatible con producto, cociente y
orientaci\'on.
\end{proposition}

\begin{proof}
De \(Z=r_-/r_+\) y \(c^{-1}=\det W=r_-r_+\) se obtienen de forma \`unica las
ra\'ices positivas. Las dos identidades constitutivas fuerzan despu\'es
\(\mu=r_-^2\) y \(\varepsilon=r_+^2\).
\end{proof}

\subsection{Determinante angular y relaci\'on electrod\'ebil de Weinberg}

El determinante de \(T\) es

\begin{equation}
\det T=q_+q_-=D_A.
\label{rm:eq:detT}
\end{equation}

El sector constitutivo del vac\'io eval\'ua \(F(T)^2\), donde
\(F(z)=(1-z^{90})/(1-z^{120})\), mientras que el sector angular
electrod\'ebil eval\'ua \(\det T\). Se trata de dos funciones de un
antecedente com\'un y no de una factorizaci\'on horizontal entre constantes.
Esta distinci\'on se desarrolla en el \cref{rm:ch:ew}.

\begin{falsifier}
El bloque queda refutado si se viola cualquiera de estas identidades exactas:
positividad de \(I-T^{120}\), igualdad
\cref{rm:eq:three-four-completion}, relaciones \cref{rm:eq:vacuum-relations} o
compatibilidad \(E_m(\mathcal V_{m+1})=\mathcal V_m\). Una coincidencia SI
no puede compensar un fallo operatorio.
\end{falsifier}

\begin{statusbox}
Operador, espectro, completaci\'on \(3\to4\), paridad, amplificaci\'on y
secciones constitutivas forman un teorema interno exacto. Su reuni\'on en un
solo operador queda establecida en esta formulaci\'on; los datos angulares y
las rectas, derivados previamente en HMT, constituyen la estructura de partida
expl\'icita.
\end{statusbox}


\section{Determinante angular y par\'ametros del sector electrod\'ebil}
\label{rm:ch:ew}

\subsection{Funciones constitutiva y determinantal del operador angular}

Recordemos el operador de dos hojas del \cref{rm:ch:vacuum},
\begin{equation}
\label{rm:eq:ew-angular-operator}
T=q_+P_+ +q_-P_-,
\qquad
q_-=\ee^{-(A_{\rm rad}-C^*_{\rm rad})},
\qquad
q_+=\ee^{-(A_{\rm rad}+C^*_{\rm rad})}.
\end{equation}
El vac\'io electromagn\'etico aplica la funci\'on
\begin{equation}
\label{rm:eq:ew-vacuum-reader}
F(T)^2,
\qquad
F(z)=\frac{1-z^{90}}{1-z^{120}},
\end{equation}
mientras el sector angular electrod\'ebil aplica el determinante
\begin{equation}
\label{rm:eq:ew-determinant-reader}
\boxed{
\det T=q_+q_-=\ee^{-2A_{\rm rad}}=:D_A.}
\end{equation}
Son dos realizaciones funcionales hermanas del antecedente com\'un
\((A,C^*,T)\). La estructura causal genera ambas desde \(T\). La
reconstrucci\'on posterior demostrada en
\cref{rm:thm:curve-F-bijection,rm:thm:curve-vacuum-tomography} permite recuperar
\[
 \det T
 =\det\!\left(F^{-1}\!\left(\sqrt{F(T)^2}\right)\right),
\]
sin convertir esa inversi\'on funcional en la flecha generadora del sector
electrod\'ebil. La unidad estructural reside en \(T\), y la reconstrucci\'on
explicita la compatibilidad entre permeabilidad, permitividad y
\(\theta_W\).

En la involuci\'on de hojas \(C^*\mapsto-C^*\), los autovalores
\(q_+,q_-\) se intercambian. Por ello \(D_A\) y la velocidad constitutiva
son pares, mientras \((\widehat\mu,\widehat\varepsilon)\) se intercambia y
\(\widehat Z\) se invierte. La relaci\'on electrod\'ebil que sigue usa el
invariante par.

\subsection{Relaci\'on exacta entre los sectores \texorpdfstring{\(W\)}{W} y \texorpdfstring{\(Z\)}{Z}}

La coordenada angular de duraci\'on--perduraci\'on es
\begin{equation}
\label{rm:eq:ew-angular-character}
\chi_{\rm dur}(n_A,s_C)
=\exp\!\left(n_AA_{\rm rad}
+\frac{s_C}{6}C^*_{\rm rad}\right).
\end{equation}
Esta coordenada no se introduce como una etiqueta aislada. Su procedencia es
\begin{equation}
\label{rm:eq:ew-signature-provenance}
\begin{aligned}
\APP&\longrightarrow\TRIT\longrightarrow\TPK
\longrightarrow\text{extensión superviviente}
\longrightarrow\text{ruta}\\
&\longrightarrow\text{registro genealógico}
\longrightarrow(n_A,s_C,\ldots)
\longrightarrow\chi_{\rm dur}.
\end{aligned}
\end{equation}
Las firmas de \(W\), \(Z\) y \(H\) son magnitudes tipadas obtenidas mediante
esa cadena; el presente cap\'itulo las adopta sin seleccionarlas a partir de
sus masas terminales. En la realizaci\'on angular establecida, los sectores
cargado y neutro satisfacen
\begin{equation}
\label{rm:eq:ew-WZ-signatures}
(n_A,s_C)_W=(94,-1),
\qquad
(n_A,s_C)_Z=(95,-1).
\end{equation}
Esta afirmaci\'on se refiere a la componente angular com\'un. Las
coordenadas adicionales de refinamiento de la ley general de masas no se
eliminan salvo que el transductor f\'isico demuestre su cancelaci\'on.

\begin{theorem}[Defecto de Weinberg en el esquema de masas f\'isicas de la carta angular]
\label{rm:thm:ew-weinberg}
Si \(W\) y \(Z\) se eval\'uan en una misma secci\'on de escala y comparten
el refinamiento que acompa\~na a \cref{rm:eq:ew-WZ-signatures}, entonces
\begin{equation}
\label{rm:eq:ew-WZ-ratio}
\boxed{
\frac{m_W^{\angle}}{m_Z^{\angle}}
=\ee^{-A_{\rm rad}}=\sqrt{D_A},
\qquad
\left(\frac{m_W^{\angle}}{m_Z^{\angle}}\right)^2=D_A.}
\end{equation}
En el esquema de masas f\'isicas,
\begin{equation}
\label{rm:eq:ew-sin2thetaW}
\boxed{
\sin^2\theta_W^{\rm OS,HMT}=1-D_A
=0.2248708807528435698111159035\ldots .}
\end{equation}
Adem\'as,
\begin{equation}
\label{rm:eq:ew-WZ-terminal-ratio}
\frac{m_W^{\angle}}{m_Z^{\angle}}
=0.8804141748331613670128885459\ldots .
\end{equation}
\end{theorem}

\begin{proof}
Por \cref{rm:eq:ew-angular-character,rm:eq:ew-WZ-signatures},
\[
\frac{\chi_{\rm dur}(94,-1)}{\chi_{\rm dur}(95,-1)}
=\exp(-A_{\rm rad});
\]
la coordenada \(C^*\) se cancela porque ambos sectores comparten
\(s_C=-1\). Al elevar al cuadrado y usar
\cref{rm:eq:ew-determinant-reader} se obtiene \(D_A\). La definici\'on
de masas f\'isicas \(s_W^2=1-m_W^2/m_Z^2\) produce
\cref{rm:eq:ew-sin2thetaW}.
\end{proof}

El contenido principal es la identidad entre dos construcciones
internos: el determinante del operador angular y la diferencia de un paso
en la realizaci\'on \(W/Z\). La promoci\'on a selector f\'isico de calibre exige
un entrelazador que lleve las hojas de \(T\) a la base neutra
\((W^3,B)\), y que demuestre qu\'e refinamientos se conservan o cancelan.

\subsection{Acoplamientos de calibre en normalizaci\'on racionalizada}

Definimos el acoplamiento adimensional de calibre
\begin{equation}
\label{rm:eq:ew-egauge}
e_g:=\sqrt{4\pi\alpha_{\HMT}}
=0.3028221208717526427662442689\ldots .
\end{equation}
No debe confundirse \(e_g\) con una secci\'on dimensional de carga
\(e_{\rm int}\): la primera es un acoplamiento en unidades naturales; la
segunda vive en la recta de carga de la Mec\'anica Dimensional.

Sea
\begin{equation}
\label{rm:eq:ew-sw-cw}
s_W^2=1-D_A,
\qquad
c_W^2=D_A.
\end{equation}
La convenci\'on electrod\'ebil racionalizada a nivel de \`arbol da
\begin{align}
g&=\frac{e_g}{s_W}
=2\sqrt{\frac{\pi\alpha_{\HMT}}{1-D_A}}
=0.6385883427334574283647003809\ldots,
\label{rm:eq:ew-g}\\
g'&=\frac{e_g}{c_W}
=2\sqrt{\frac{\pi\alpha_{\HMT}}{D_A}}
=0.3439541633108502429362319257\ldots,
\label{rm:eq:ew-gprime}\\
\frac{g'}{g}
&=\sqrt{D_A^{-1}-1}
=0.5386164141966093594996610933\ldots .
\label{rm:eq:ew-gprime-over-g}
\end{align}

\begin{proposition}[Relaci\'on custodial de \`arbol]
\label{rm:prop:ew-rho}
En la misma interfaz,
\begin{equation}
\label{rm:eq:ew-rho}
\rho_{\rm EW}^{(0)}
:=\frac{(m_W^{\angle})^2}
{(m_Z^{\angle})^2c_W^2}=1.
\end{equation}
\end{proposition}

\begin{proof}
Por \cref{rm:eq:ew-WZ-ratio,rm:eq:ew-sw-cw}, tanto el numerador relativo como
\(c_W^2\) son \(D_A\).
\end{proof}

Las ecuaciones \cref{rm:eq:ew-g,rm:eq:ew-gprime} son composiciones exactas una
vez declarados el esquema de masas f\'isicas, la racionalizaci\'on y la escala. No
afirman que un cambio de esquema de renormalizaci\'on deje inalterados sus
valores num\'ericos.

\subsection{Constante de Fermi en aproximaci\'on de \`arbol}

En unidades naturales, la integraci\'on del propagador cargado a baja
energ\'ia produce
\begin{equation}
\label{rm:eq:ew-GF-def}
\frac{G_F^{(0)}}{\sqrt2}=\frac{g^2}{8m_W^2}.
\end{equation}
Sustituyendo \cref{rm:eq:ew-g} se obtiene la composici\'on
\begin{equation}
\label{rm:eq:ew-GF}
\boxed{
G_F^{(0)}
=\frac{\pi\alpha_{\HMT}}
{\sqrt2(1-D_A)m_W^2}.}
\end{equation}
Al adoptar la masa obtenida previamente
\begin{equation}
\label{rm:eq:ew-W-output}
m_W^{\HMT}=80.381592550221\ldots\ \mathrm{GeV},
\end{equation}
la evaluaci\'on terminal es
\begin{equation}
\label{rm:eq:ew-GF-value}
\boxed{
G_F^{(0)}
=1.1157162817659203186621784\times10^{-5}\ \mathrm{GeV}^{-2}.}
\end{equation}

La diferencia respecto de una constante de Fermi renormalizada no se usa
para elegir ning\'un coeficiente. En la convenci\'on
\begin{equation}
\label{rm:eq:ew-Delta-r}
G_F
=\frac{\pi\alpha_{\HMT}}
{\sqrt2(1-D_A)m_W^2(1-\Delta r)},
\end{equation}
la obligaci\'on restante es derivar \(\Delta r\) desde el sector interno de
bucles, estados y escala. Abrir el valor observado de \(G_F\) para definir
\(\Delta r\) ser\'ia una reconstrucci\'on calibrada, no una predicci\'on.

\subsection{Dependencia orientada del acoplamiento de Higgs}

En la carta angular, las rutas pertinentes son
\begin{equation}
\label{rm:eq:ew-HW-signatures}
(n_A,s_C)_H=(97,9),
\qquad
(n_A,s_C)_W=(94,-1).
\end{equation}
Por tanto,
\begin{equation}
\label{rm:eq:ew-HW-ratio}
\boxed{
\frac{m_H^{\angle}}{m_W^{\angle}}
=\frac{\chi_{\rm dur}(97,9)}{\chi_{\rm dur}(94,-1)}
=\exp\!\left(3A_{\rm rad}+\frac53C^*_{\rm rad}\right)
=1.55872087212325548564\ldots .}
\end{equation}
Aqu\'i \(C^*\) no se cancela: el sector escalar conserva la orientaci\'on
que el determinante \(D_A=\ee^{-2A_{\rm rad}}\) no distingue.

Tomemos la convenci\'on del potencial
\begin{equation}
\label{rm:eq:ew-Higgs-potential}
V(H)=-\mu_H^2H^\dagger H
+\lambda_H(H^\dagger H)^2,
\end{equation}
para la cual
\begin{equation}
\label{rm:eq:ew-Higgs-tree-relations}
m_H^2=2\lambda_Hv^2,
\qquad
m_W^2=\frac{g^2v^2}{4}.
\end{equation}

\begin{corollary}[Autoacoplamiento orientado de Higgs]
\label{rm:cor:ew-Higgs-lambda}
En la interfaz de \`arbol declarada,
\begin{equation}
\label{rm:eq:ew-Higgs-lambda}
\boxed{
\lambda_H
=\frac{\pi\alpha_{\HMT}}{2(1-D_A)}
\exp\!\left(6A_{\rm rad}+\frac{10}{3}C^*_{\rm rad}\right)
=0.1238479115482466804662\ldots .}
\end{equation}
Adem\'as,
\begin{equation}
\label{rm:eq:ew-GF-Higgs}
G_F^{(0)}m_H^2=\sqrt2\,\lambda_H.
\end{equation}
\end{corollary}

\begin{proof}
Dividir las dos ecuaciones de
\cref{rm:eq:ew-Higgs-tree-relations} da
\(m_H^2/m_W^2=8\lambda_H/g^2\). Con
\(g^2=4\pi\alpha_{\HMT}/(1-D_A)\) y el cuadrado de
\cref{rm:eq:ew-HW-ratio} se obtiene \cref{rm:eq:ew-Higgs-lambda}. Por otra parte,
\(G_F^{(0)}=1/(\sqrt2v^2)\) y
\(m_H^2=2\lambda_Hv^2\) implican
\cref{rm:eq:ew-GF-Higgs}.
\end{proof}

El valor de \(\lambda_H\) es anterior a la incorporaci\'on de correcciones
radiativas y depende de
la convenci\'on y la escala declaradas. Su inter\'es estructural es que
separa dos memorias: el sector com\'un de calibre emplea el invariante par \(D_A\),
mientras la ruta de Higgs retiene expl\'icitamente \(C^*\).

\subsection{Clausura determinantal del sector electrodébil: mezcla, acoplamientos y espectro constitutivo}

La cadena completa de este cap\'itulo es
\begin{equation}
\label{rm:eq:ew-causal-diagram}
\boxed{
\begin{aligned}
(A,C^*)&\longrightarrow T
\longrightarrow
\begin{cases}
F(T)^2&\longrightarrow
(\widehat\mu,\widehat\varepsilon,\widehat Z,\widehat c),\\
\det T=D_A&\longrightarrow
(s_W^2,c_W^2),
\end{cases}\\
(D_A,\alpha_{\HMT})&\longrightarrow(e_g,g,g'),\\
(D_A,\alpha_{\HMT},m_W)&\longrightarrow G_F^{(0)},\\
(A,C^*,D_A,\alpha_{\HMT})&\longrightarrow
\left(\frac{m_H}{m_W},\lambda_H\right).
\end{aligned}}
\end{equation}
La cifra aparece siempre al final. Ni \(\sin^2\theta_W\), ni
\(G_F\), ni \(\lambda_H\) se introducen como valores objetivo para seleccionar
\(A,C^*\) o las firmas de ruta.

\begin{keyresult}
El mismo operador angular determina el vac\'io mediante una funci\'on racional
de sus hojas y el defecto electrod\'ebil mediante su determinante. La
diferencia de un paso entre \(W\) y \(Z\) convierte ese determinante en
\(\sin^2\theta_W^{\rm OS}=1-D_A\); la ruta de Higgs, en cambio, conserva
el contra\'angulo y conserva la orientaci\'on eliminada por el determinante.
\end{keyresult}

\begin{falsifier}
La relaci\'on de Weinberg queda refutada en su dominio si las firmas
angulares \(W/Z\) no difieren s\'olo por \(n_A:94\to95\), si una coordenada
de refinamiento no se cancela en la interfaz declarada o si
\((m_W^{\angle}/m_Z^{\angle})^2\neq D_A\). La promoci\'on de calibre queda
bloqueada si no existe un entrelazador con la base \((W^3,B)\). Las cifras
de \(g,g',G_F,\lambda_H\) cambian si se altera esquema o escala sin
registrarlo. Usar \(G_F\) o \(m_H\) observados para seleccionar
\(\Delta r\), las firmas o \(C^*\) introduce retroalimentaci\'on por el
valor objetivo e invalida la derivaci\'on. Confundir \(e_g\) con la carga
dimensional constituye un error de tipo.
\end{falsifier}

\begin{statusbox}
Determinante angular, cociente \(W/Z\) y defecto complementario de Weinberg:
\status{Teorema interno exacto en la carta angular condicionada}. Acoplamientos \(g,g'\),
\(\rho_{\rm EW}^{(0)}\), Fermi y Higgs: \status{Composici\'on exacta en la
interfaz electrod\'ebil de \`arbol}. El entrelazador de base de calibre y la
correcci\'on interna \(\Delta r\): \status{Frontera abierta con criterio de refutación}.
Fermi y Higgs no son constantes independientes seleccionadas sin valores
objetivo: son composiciones exactas que emplean las firmas, la magnitud
\(m_W\) y la
convenci\'on de \`arbol declarada. Los valores convencionales posteriores
son \status{Contraste metrol\'ogico externo}.
\end{statusbox}


\paragraph{Del determinante electrodébil al operador constitutivo positivo del vacío y su completación espectral \(3\to4\)} Una vez fijado el soporte constitutivo, la evolución cuántica se reconstruye como suma e integral de genealogías, con reducción efectiva y publicación probabilística tipadas.
